% Generated draft: requires mathematical and source review.
\documentclass[11pt]{article}
\usepackage[margin=25mm]{geometry}
\usepackage{fontspec}
\setmainfont{Latin Modern Roman}
\usepackage{amsmath,amssymb,mathtools,unicode-math}
\setmathfont{Latin Modern Math}
\usepackage{enumitem,needspace,xurl,hyperref}
\setlength{\parindent}{0pt}
\setlength{\parskip}{5pt}
\setlength{\emergencystretch}{3em}
\newcommand{\N}{\mathbb{N}}
\newcommand{\Z}{\mathbb{Z}}
\newcommand{\Q}{\mathbb{Q}}
\newcommand{\R}{\mathbb{R}}
\newcommand{\C}{\mathbb{C}}
\newcommand{\F}{\mathbb{F}}
\newcommand{\A}{\mathbb{A}}
\newcommand{\PP}{\mathbb{P}}
\newcommand{\E}{\mathbb{E}}
\newcommand{\eps}{\varepsilon}
\newcommand{\dd}{\,\mathrm{d}}
\begin{document}
\section*{UnsolvedMath ID: 1510}
\textbf{Source problem code:} catalogue entry
\subsection{Definitions and mathematical statement}
A complex number $\alpha$ is transcendental over $\Q$ if $P(\alpha)\ne0$ for every nonzero polynomial $P\in\Q[T]$. Two complex numbers $a,b$ are algebraically independent over $\Q$ if
\[
\forall P\in\Q[X,Y]\setminus\{0\},\qquad P(a,b)\ne0.
\]
Let $e=\exp(1)$ and let $\pi$ be the usual positive circle constant. The displayed source question asks whether $e$ and $\pi$ are algebraically independent, and whether each of
\[
e+\pi,\qquad e\pi,\qquad \pi^e
\]
is transcendental. Its background additionally names $\pi^\pi$. Since the bases are positive real numbers, use the real-power convention $\pi^a=\exp(a\log\pi)$ for $a=e$ or $a=\pi$, with the real logarithm. The phrase ``etc.'' does not identify any further testable list of constants.

For precision, the full standard Schanuel conjecture named in the title is a broader assertion: for every integer $n\geq1$ and every tuple $z_1,\ldots,z_n\in\C$ that is linearly independent over $\Q$,
\[
\operatorname{trdeg}_{\Q}\Q(z_1,\ldots,z_n,\exp z_1,\ldots,\exp z_n)\geq n.
\]
Linear independence here means that $\sum_{j=1}^n q_jz_j=0$ with $q_j\in\Q$ forces all $q_j=0$. Transcendence degree is the size of a transcendence basis, equivalently the largest size of an algebraically independent subset for this finitely generated field. The displayed Schanuel statement is the standard conjecture identified by the title, whereas the stored question asks about particular constants and does not assert that proving one of these special cases proves the full conjecture.

The supplied background says that Schanuel's conjecture would imply algebraic independence of $e$ and $\pi$ and the stated transcendence consequences $e+\pi$, $e\pi$, and $\pi^\pi$. The main sentence separately asks about $\pi^e$; it is kept distinct from $\pi^\pi$. The supplied triage dated 2026-08-17 labels the topic open, says the two-variable Schanuel regime is already difficult and would imply algebraic independence of $e$ and $\pi$, and cites Lindemann--Weierstrass only for weaker individual transcendence facts. It reports no verified unconditional algebraic-independence theorem for $e$ and $\pi$ or resolution of the concrete open questions listed in the record. This is a dated source assessment, not an independent verification or a claim that all conceivable expressions involving $e$ and $\pi$ are open.
\subsection{Short English statement}
Are $e$ and $\pi$ algebraically independent, and are $e+\pi$, $e\pi$, $\pi^e$, and the background's separate example $\pi^\pi$ transcendental? The title points to the broader Schanuel conjecture, whose precise universal statement is distinguished from these concrete questions.
\subsection{Sources}
\begin{itemize}
\item D. Bertrand, G. Binyamini and M. Kowalski, ``Exponential sums equations and tropical geometry,'' Selecta Mathematica 29 (2023), as supplied: \url{https://doi.org/10.1007/s00029-023-00853-y}. The supplied record cites this for the difficult two-variable case and its implication for $e$ and $\pi$.
\item Local registry data/unsolvedmath/problems.json, permanent catalogue entry, source code catalogue entry, ``Schanuel's Conjecture''; statement, background, and literature triage dated 2026-08-17. The background mentions Lindemann--Weierstrass without a separate bibliographic reference. Dataset attribution: Ulam AI and the UnsolvedMath Contributors.
\end{itemize}
\end{document}
